% SPDX-License-Identifier: Apache-2.0
\section{The Configuration Flash After the Bitstream}
\label{sec:x-cfgflash}

The master of Section~\ref{sec:x-spi} and the window of
Section~\ref{sec:x-flashwin} each reach a flash on four wires of their
own. The flash worth reaching on the board is the one that holds the
bitstream, and its clock is not a wire of anyone's. On a 7-series part
it is the dedicated CCLK pin, which the configuration logic drives
while the bitstream loads and which user logic reaches afterwards only
through the \code{STARTUPE2} primitive's \code{USRCCLKO} input.

The primitive brings two rules with it. Its \code{EOS} output says
when configuration is over, and until then the flash is not the user's.
The first three clocks on \code{USRCCLKO} after \code{EOS} never reach
the pin, since the part spends them moving the pin over, so a command
whose clocks start there reaches the chip three bits short.

\code{CfgFlash} is the answer, in two pieces. \code{Startup} is the
primitive as a foreign module: synthesised, it is \code{STARTUPE2};
simulated, in Rust, in Verilog and in VHDL alike, it is a model that
raises \code{eos} after a few cycles and swallows the three clocks, and
it brings out the pin as \code{cclk} so that a model of the chip can
be clocked by it. \code{FlashPins} stands between the chip and the two
parts. It waits for \code{eos}, spends four clocks with the chip not
selected, and holds the window in reset until then. After that it
gives the wires to whichever part selects the chip first, the window
before the master when both ask at once, until that part lets go.

\code{FlashPins} also refuses write enable, \code{0x06}, from the
master. Every command that programs or erases a flash, or writes its
registers, is ignored by the chip unless write enable came first, so
that one command is the door to all of them. With \code{GUARD} set,
the pins let the chip go on the seventh bit of a write enable, and the
chip discards a command shorter than eight bits. Nothing the core runs
can then change the flash that holds the bitstream.

Every output of the pins is a register, and the clock passes the
primitive as well, so the chip's answer comes back two cycles after
the edge that asked for it. The run clocks both parts with a half bit
of four cycles for that reason.

The run asks the window for a word in its first cycle, before
configuration is over. The read waits and is answered once the chip is
ready. The master then asks the chip who it is, which is the first
command after the swallowed clocks and arrives whole. It sends write
enable, which the pins refuse, and the window reads a second word to
show the chip still answers. At the end the chip model lists the
commands it took whole: the two reads and the identity, and never
write enable.

\begin{figure}[htbp]
\centering
\input{cfgflash_timing.tex}
\caption{The start: \code{phase} leaves its wait at the end of startup,
the pins turn \code{usrcclko} four times with the chip not selected,
and only the fourth turn reaches \code{cclk}. Then the window is let
out of reset and its read selects the chip.}
\label{fig:x-cfgflash}
\end{figure}

\lstinputlisting[caption={\code{ex\_cfgflash.rs}. Run with \code{bazel run //lib/examples:ex\_cfgflash}.},label={lst:x-cfgflash}]{lib/examples/ex_cfgflash.rs}

\lstinputlisting[style=ascii,caption={What \code{ex\_cfgflash} printed when the build ran it: the two words, the identity, the refused command, and what the chip took whole.},label={lst:x-cfgflash-out}]{out_cfgflash.txt}
